\documentclass[11pt,a4paper]{article}
\usepackage[T1]{fontenc}
\usepackage{lmodern}
\usepackage[margin=25mm]{geometry}
\usepackage{amsmath,amssymb,amsthm,mathtools}
\usepackage{mathrsfs}
\usepackage{booktabs,longtable,array}
\usepackage{xcolor,microtype}
\usepackage{enumitem}
\usepackage{listings}
\usepackage{xurl}
\usepackage{hyperref}
\hypersetup{colorlinks=true,linkcolor=blue!45!black,citecolor=blue!45!black,urlcolor=blue!45!black,
pdftitle={Additive coding and smaller coefficients for Diophantine consistency equations},
pdfauthor={AI-assisted research report prepared for ProveIt},
pdfsubject={Explicit Diophantine presentations of Con(ZFC)},
pdfkeywords={Diophantine equations, Hilbert tenth problem, ZFC consistency, Bose-Chowla, additive coding}}
\setlength{\emergencystretch}{3em}
\setlist{itemsep=3pt,topsep=5pt}
\newtheorem{theorem}{Theorem}[section]
\newtheorem{lemma}[theorem]{Lemma}
\newtheorem{proposition}[theorem]{Proposition}
\newtheorem{corollary}[theorem]{Corollary}
\theoremstyle{definition}
\newtheorem{definition}[theorem]{Definition}
\newtheorem{remark}[theorem]{Remark}
\newcommand{\N}{\mathbb N}
\newcommand{\Z}{\mathbb Z}
\newcommand{\Q}{\mathbb Q}
\newcommand{\F}{\mathbb F}
\newcommand{\Con}{\operatorname{Con}}
\newcommand{\pc}{\operatorname{popcount}}
\newcommand{\code}[1]{\ifmmode\text{\nolinkurl{#1}}\else\nolinkurl{#1}\fi}
\newcommand{\B}{\mathcal B}
\newcommand{\capc}{\mathscr C}
\DeclareMathOperator{\Sq}{Sq}
\title{\textbf{Additive coding and smaller coefficients\\for Diophantine consistency equations}\\[4pt]
\large An explicit synthesis for $\Con(\mathrm{ZFC})$}
\author{Research report prepared for ProveIt\\
\small AI-assisted mathematical synthesis, construction, and independent audit}
\date{23 October 2026}
\begin{document}
\maketitle
\begin{abstract}
We construct a completely specified ordinary polynomial in ten unrestricted
integer unknowns whose unsolvability is equivalent, provably in ZFC, to
$\Con(\mathrm{ZFC})$ for the repository's proof predicate. Its degree is
exactly $107\,286\,763\,576$, compared with
$126496\cdot5^{31}+27064$ for the ten-integer construction in the pinned
repository. The degree decreases by a factor of approximately
$5.49\cdot10^{15}$. The arithmetic circuit grows from 1,442 to 2,204 binary
operations. This is a degree improvement at a fixed number of unknowns.

The new step replaces exponentially spaced positions $5^j$ by a finite
additive $B_4$ code. Quartic coefficient recovery needs only distinct
sums of at most four positions. A translated Bose--Chowla set supplies
32 positions, including the homogenizing zero, with maximum $942302$.
We prove the exact coefficient identity, bounds, mask decoding, and
the complete source interface required by Sun's ten-integer reduction.

Independently, a native compiler refinement encodes every edge label
by a nonzero word, making separate occupancy bits redundant. Together
with an audited support layout, this reduces the native base from the
previous smallest bound $2^{1012115}$ to $2^{683195}$. The four fixed
coefficient files occupy 342,269 raw bytes. We bind this tuple to the
existing P67 equation and 59-operation system. We also deliver a
fully composed ten-integer variant, regenerating all source coefficients
and bounds from the new tuple. It retains the same degree and circuit
size and reduces its five raw fixed payloads from 1,066,835 to
685,209 bytes. A separate comparison instance preserves the original
tuple and isolates the additive-code improvement.

The archive includes the report source and PDF, complete circuits,
actual fixed coefficients, reproducible generators, finite
certificates, independent reviews, and provenance. No global
minimality, new universal variable-count theorem, consistency proof,
or new end-to-end Lean/Rocq certificate is claimed.
\end{abstract}
\clearpage
\tableofcontents
\newpage

\section{The results and what they improve}
\label{sec:result}
All natural numbers include zero. A positive variable is strictly
positive. An unrestricted integer variable may have any sign.
These domains are part of each theorem.

Let $\operatorname{Ref}_{\mathrm{ZFC}}$ denote the existence of a
finite refutation in the proof calculus used by the pinned ProveIt
development. Thus $\Con(\mathrm{ZFC})$ abbreviates
$\neg\operatorname{Ref}_{\mathrm{ZFC}}$ for that calculus.

\begin{theorem}[The comparison ten-integer equation]\label{thm:main}
The comparison circuit \code{ten_integer/generated/b4_signed10.json}
defines a polynomial
\[
 F_{\mathrm{B4}}(f,g,h,k,l,w,x,y,z,m)\in
 \Z[f,g,h,k,l,w,x,y,z,m]
\]
such that
\[
 \mathrm{ZFC}\ \vdash\
 \bigl(\exists\boldsymbol u\in\Z^{10}:F_{\mathrm{B4}}(\boldsymbol u)=0\bigr)
 \longleftrightarrow \operatorname{Ref}_{\mathrm{ZFC}}.
\]
Consequently,
\[
 \boxed{\quad
 \mathrm{ZFC}\ \vdash\
 \bigl(\nexists\boldsymbol u\in\Z^{10}:F_{\mathrm{B4}}(\boldsymbol u)=0\bigr)
 \longleftrightarrow\Con(\mathrm{ZFC}).
 \quad}
\]
The exact degree is $107286763576$. The supplied evaluator uses
485 additions, 25 subtractions, and 1,694 multiplications:
2,204 binary arithmetic operations in total.
\end{theorem}

The degree decrease is
\[
\frac{126496\cdot5^{31}+27064}{107286763576}
=
\frac{589042901992797851562527064}{107286763576}
\approx5.4903595034\cdot10^{15}.
\]
There are 762 more arithmetic operations than in the previous
ten-integer circuit, an increase of about $52.8\%$.
The new method improves degree dramatically while retaining the
same ten unrestricted unknowns and a compact, finite expression.

The recommended complete instance is
\code{ten_integer_composed/generated/b4_signed10.json}.
Corollary~\ref{composed:main} proves the same equivalence, degree,
and operation count for this variant, with all fixed coefficients
and the source-base bound recomputed from the smaller tuple below.
The comparison circuit remains included to separate the two
mathematical changes.

The second result concerns a different measure:
\begin{theorem}[A smaller complete native tuple]\label{thm:tuple}
The nonzero-edge compiler in \code{coefficient_backend/} supplies
an admissible fixed blank-K215 tuple
$(B-1,K,M_C,M_F^{\rm src})$ with
\[
 B=2^{683195},\qquad
 \max\{\text{port bit lengths}\}=688546.
\]
Its four materialized unsigned integer files contain 342,269 raw
bytes. Substitution of this entire tuple in the inherited P67
equation and retained 59-operation system preserves their
equivalence to blank K215 halting, and hence their unsolvability
is equivalent in ZFC to $\Con(\mathrm{ZFC})$.
\end{theorem}
The source base exponent decreases by $32.50\%$ compared with
package 36's $1012115$; raw bytes decrease by $32.45\%$ from
506,664. These percentages compare complete tuples. The proof
does not splice the smallest individual ports from different
constructions.

\subsection{There is no single smallest measure}
Variable count, degree, number of arithmetic operations, coefficient
size, and proof complexity are distinct. The present work supplies
two new points in that comparison. It does not dominate the old
67-operation equation or the nine-natural-variable constructions
on every measure.

\begin{center}
\small
\begin{tabular}{@{}p{44mm}p{23mm}p{43mm}r@{}}
\toprule
Construction & Domain & Degree & Operations\\
\midrule
Pinned package 31 & $10$ integers &
$126496\cdot5^{31}+27064$ & 1442\\
This report, main circuit & $10$ integers &
$107286763576$ & 2204\\
Pinned P67; new tuple here & $15$ positive &
$155$ & 67\\
Pinned quartic, package 31 & $31$ positive &
$4$ & 125\\
Pinned package 34 & $12$ positive &
$1022$ & 124\\
Pinned package 37 & $9$ positive &
$92324\cdot5^{32}+296$ & 1500\\
\bottomrule
\end{tabular}
\end{center}
The pinned values are the audited collection's stated values,
supported by their source formulas. The new circuits and
coefficient metrics are recomputed in this delivery.
Nine-natural or nine-positive results are not ten-integer results
with a different notation.

\section{Repository analysis and the inherited boundary}
\label{sec:repo}
The source revision is
\[
 \code{22e75dd6b87b1b199bb08c778d7cf075ae5aae08}.
\]
Its collection contains wave 5, packages 31--37, with a subsequent
audit correcting several comparisons. Earlier waves remain
recoverable from history. In particular, the relevant baseline
already has ten unrestricted integers via Sun's 2026 preprint
\cite{Sun2026}; improving to fourteen or twelve unrestricted
integers would not advance that frontier.

\subsection{What was combined}
Package 31 provides the explicit ten-integer outer construction,
its quintic divisibility merger, the coprime positivity merger,
and economical separating weights. It also supplies a complete
fixed coefficient tuple and the 31-variable quartic. These are
retained in the main result.

The new additive code replaces the source polynomial's
exponentially spaced digit positions. It is not an optimization
of a few gates in the final norm. It changes the exponent geometry
at the stage that determines the enormous degree.

The coefficient result combines the Wang943 native compiler,
the actual-rank control method refined in package 36, a new
nonzero-word encoding of edge labels, and an exact support
layout. It improves explicit coefficient size while preserving
the simple P67 and 59-operation arithmetic skeletons.

The inspections of packages 34 and 37 also identify two relevant
boundaries: their shifted twelve-positive-variable norms are
different polynomials whose positive completeness requires a
proof, and their small signed Jones forms rely on additional
cofinality arguments. Neither of those delicate transformations
is needed for the principal ten-integer result here.

\subsection{The fixed parent}
The positive variables, in their authoritative order, are
\[
\begin{split}
 &(J,F,\alpha,z,f,h,i,o,j,s,w,\tau,\eta,\zeta,y,\gamma,\\
 &\hspace{12mm}N,E,P,c,a,X,D,C_2,L,F_2,V,g,I,B_{\rm aux},\Delta).
\end{split}
\]
The four fixed ports are $(\beta,K_0,M_C,M_F)$, with $M_F$
already the shifted mask. Set
\[
 q=\beta J+1,\qquad k=\eta+\zeta,\qquad
 H=4a+3,\qquad C=q-F-\alpha.
\]
Appendix~\ref{app:parent} gives all 22 quadratic residuals $r_j$.
For the nonnegative coding domain define
\begin{equation}\label{eq:parent}
 P_0(n_1,\ldots,n_{31})=
 \sum_{j=1}^{22}r_j(n_1+1,\ldots,n_{31}+1)^2.
\end{equation}
This is a quartic. Its constant term is strictly positive because
at the all-one parent assignment the residual $E-X(a-E)$ is $1$.
Every zero therefore has at least one nonzero coordinate.

The comparison tuple is package 31's complete 121-by-8801
tuple, whose native base is $2^{1064921}$. Its four materialized
ports have bit lengths
\[
 1064921,\quad1070004,\quad1064921,\quad1064922.
\]
The compressed files and \code{ten_integer/data/INSTANCE.json}
fix them exactly. Historical metadata in a byte-preserved parent
file does not override this active binding.
Section~\ref{composed:section} binds the same parent equations to
the smaller tuple and recomputes all following coefficient bounds.

For
\[
 P_0(\boldsymbol n)=\sum_{|\alpha|\le4}p_\alpha\boldsymbol n^\alpha,
 \qquad
 c_\alpha=(4-|\alpha|)!\prod_{i=1}^{31}\alpha_i!\,p_\alpha,
\]
the exact reconstruction gives 402 nonzero records and proves
\[
 |c_\alpha|+1<2^{4269855}.
\]
The source radix multiplier is the actual fixed integer
\[
 \beta_S=2^{4269882},
\]
which is also materialized and hashed. Since
$48\cdot33^4=56924208<2^{27}$, it satisfies
\[
 \beta_S>48\cdot33^4\max_\alpha|p_\alpha|.
\]
All coefficient records are finite integer expressions in the
four fixed ports. There is no unspecified recursively enumerable
index or existential coefficient.

\subsection{The endpoint being inherited}
The pinned development provides
\begin{equation}\label{eq:inherited}
 \exists\boldsymbol n\in\N^{31}:P_0(\boldsymbol n)=0
 \quad\Longleftrightarrow\quad
 \text{K215 halts on the blank tape}
 \quad\Longleftrightarrow\quad
 \operatorname{Ref}_{\mathrm{ZFC}}.
\end{equation}
The formal machine theorem is
\code{KEndpoint215.blank_halts_iff_zfc_refutation}, whose statement
uses \code{BlankHalts machine215} and
\code{BProv ZFCax_s [] Form.fBot}.
The native compiler and quartic adapters remain ordinary
mathematical developments. Their particular proof files and
the coefficient adapter are cited and retained with provenance.

We use the mathematical equivalences represented by those
developments. We do not infer a ZFC theorem from Lean's soundness,
and we do not claim a new build of the formal machine endpoint.

\section{Additive coding of the quartic source}
\label{b4:section}

The main new construction replaces the positional exponents in the
ten-integer equation of package~31 by an additive code.  The source
quartic, its complete coefficient tuple, and the outer Lucas and
relation-combining formulas are retained.  In particular, the
ten-integer theorem remains Sun's theorem \cite{Sun2026}; the result
below is an explicit improvement of its source specialization
\cite{Package31}.  Bose--Chowla codes themselves are classical
\cite{BoseChowla}.

\begin{theorem}[The explicit ten-integer instance]
\label{b4:main}
The circuit \code{ten_integer/generated/b4_signed10.json} defines a
polynomial
\[
 F_{B_4}\in\mathbb Z[f,g,h,k,l,w,x,y,z,m]
\]
of exact total degree
\[
 \boxed{\deg F_{B_4}=107\,286\,763\,576.}
\]
Its supplied division-free circuit has $2\,204$ binary additions,
subtractions, and multiplications.  With the fixed parent quartic and
compiler tuple specified in this report, ordinary mathematics
formalizable in ZFC proves
\[
 \exists(f,g,h,k,l,w,x,y,z,m)\in\mathbb Z^{10}\;(F_{B_4}=0)
 \quad\Longleftrightarrow\quad
 \neg\operatorname{Con}(\mathrm{ZFC}).
\]
Consequently absence of integer zeros is equivalent, in ZFC, to
$\operatorname{Con}(\mathrm{ZFC})$.
\end{theorem}

The proof imports the same K215 calibration, native compiler, positive
quartic reduction, and Sun bridge as the pinned construction.  It gives
a new proof of the entire source-coding interface.  Exact finite checks
supplement that proof; they are not an end-to-end formal verification of
the imported results.  Sun's ten-integer bridge is cited in its 2026
preprint version \cite{Sun2026}.  No consistency proof, globally minimal
equation, or new general ten-variable theorem is claimed.

\subsection{The fixed source and its coefficient convention}
\label{b4:source}

Let $r_1,\ldots,r_{22}$ be the fixed positive-parent residuals specified
in this report, in their literal order, and put
\[
 P_0(n_1,\ldots,n_{31})
 =\sum_{j=1}^{22}r_j(n_1+1,\ldots,n_{31}+1)^2.
\]
The inherited endpoint states
\[
 \exists\mathbf n\in\mathbb N^{31}\;P_0(\mathbf n)=0
 \quad\Longleftrightarrow\quad
 \text{ZFC has a finite syntactic refutation}.
 \tag{B4.1}\label{b4:endpoint}
\]
Here $\mathbb N$ includes zero.  The positive-to-natural shift is part
of the definition and cannot be omitted.

The four fixed coefficient ports are the complete Wang943 tuple of
package~31, with port order $(\beta,K_0,M_C,M_F)$ and bit lengths
\[
 (1064921,\ 1070004,\ 1064921,\ 1064922).
\]
The high mask $M_F$ is already shifted.  None of these ports is
existentially quantified, and no port is substituted from another
tuple.  The exact parent and compressed integer payloads are included
in \code{ten_integer/data/}.

The source has degree four and $P_0(\mathbf0)>0$.  For the latter fact,
the fifth residual $E-X(a-E)$ equals $1$ at the all-one parent
assignment.  Write
\begin{align}
 P_0(\mathbf n)&=\sum_{|\alpha|\le4}p_\alpha\mathbf n^\alpha,
 &
 c_\alpha&=(4-|\alpha|)!\prod_{i=1}^{31}\alpha_i!\,p_\alpha,
 \label{b4:coeff}\\
 H&=\max_\alpha|p_\alpha|,
 &
 \beta_S&=2^{4269882}.
 \label{b4:beta}
\end{align}
The included coefficient certificate proves
\[
 \beta_S>2\cdot24\cdot33^4H.
 \tag{B4.2}\label{b4:bound}
\]
Indeed it bounds every weighted coefficient by
$|c_\alpha|+1<2^{4269855}$, and
$2\cdot24\cdot33^4=56924208<2^{27}$.  Since the integer factorial
multiplier in \eqref{b4:coeff} is positive,
$|p_\alpha|\le|c_\alpha|$.  Thus \eqref{b4:bound} follows without
requiring decimal expansion of the enormous coefficients.

The source has $402$ nonzero weighted coefficient records.  Each is
an explicit small polynomial in the four fixed ports.  The verifier
reconstructs every record from the $22$ shifted residual squares.
The leading zero retained in each serialized exponent vector is a
historical coefficient-interface convention: it is not an input digit,
an unknown, or an extra source witness.

\subsection{What the exponent code must distinguish}
\label{b4:injectivity}

\begin{definition}
A list $0=w_0<w_1<\cdots<w_{31}$ is a quartic additive code if all
sums of four list elements, allowing repetition and disregarding
order, are distinct.
\end{definition}

Padding a multi-index with $4-|\alpha|$ copies of $w_0$ shows that this
condition is equivalent to injectivity of
\[
 \alpha\longmapsto\alpha\cdot\mathbf w
 :=\sum_{i=1}^{31}\alpha_iw_i
 \qquad (|\alpha|\le4).
 \tag{B4.3}\label{b4:injective}
\]
The positional weights $5^i$ satisfy this property.  They are not
needed for either the witness-digit mask or coefficient extraction.
An additive code gives much smaller exponents.

The code used here is
\[
\begin{gathered}
0,\ 42609,\ 49910,\ 91724,\ 155840,\ 213052,\ 290241,\ 337708,\\
340502,\ 351998,\ 388348,\ 417182,\ 480443,\ 513128,\ 534148,\\
563475,\ 579595,\ 636647,\ 639228,\ 692126,\ 725139,\ 732225,\\
748536,\ 754497,\ 817940,\ 820267,\ 829985,\ 858576,\ 870324,\\
879730,\ 907673,\ 942302.
\end{gathered}
\]
The $i$th positive entry is assigned to the $i$th coordinate of the
literal parent ordering.  Throughout this section set
\[
 T=942302,\qquad L_c=4T=3769208,\qquad M=2^{20}-1=1048575.
 \tag{B4.4}\label{b4:span}
\]

\begin{lemma}[Finite code certificate]
\label{b4:code}
The displayed list is a quartic additive code.  In fact, its four-term
multiset sums are distinct modulo $M$.
\end{lemma}

\begin{proof}
There are $\binom{35}{4}=52360$ relevant multisets.  The included
short checker enumerates all of them and verifies distinct residues
modulo $M$.  This provides a direct finite certificate independent of
the construction history.

For a mathematical construction, form the binary quotient by
$X^{20}+X^3+1$ and let $\theta$ be the residue of $X$.  Exact modular
computations give $\theta^M=1$ and
\[
 \theta^{M/p}\ne1\qquad(p=3,5,11,31,41),
\]
the complete set of prime divisors of $M$.  Hence $\theta$ has order
$M$.  Its $M$ distinct powers and zero exhaust the $2^{20}$ residues.
Every nonzero residue is therefore a unit, proving that the quotient
is a field, rather than assuming irreducibility in this step.

Let $\mathbb F_{32}$ be its $32$-element subfield and choose the
$32$ logarithms $a_c$ satisfying
\[
 \theta^{a_c}=\theta+c\qquad(c\in\mathbb F_{32}).
\]
The element $\theta$ has degree four over $\mathbb F_{32}$.  If two
sums of four logarithms agree modulo $M$, the corresponding products
of four factors $\theta+c$ agree.  Subtracting their monic
degree-four polynomials leaves a polynomial of degree at most three
over $\mathbb F_{32}$.  It must vanish identically, so unique
factorization identifies the two multisets of factors.  This is the
Bose--Chowla argument \cite{BoseChowla}.

Translate the logarithms by $-535448$ modulo $M$ and take their
representatives in $[0,M)$.  The sorted list is the displayed code.
Translation preserves equality or inequality of four-term sums
modulo $M$.  The chosen origin follows the largest cyclic gap of
the original logarithms and reduces the maximum representative to $T$.
The generator reproduces each finite-field calculation, all logarithms,
the translation, and the exhaustive multiset certificate.
\end{proof}

\subsection{The coefficient identity}
\label{b4:polarization}

Define formal polynomials
\begin{align}
 C_{\mathbf n}(t)&=\left(1+\sum_{i=1}^{31}n_it^{w_i}\right)^4,
 &
 D_c(t)&=\sum_{|\alpha|\le4}c_\alpha t^{L_c-\alpha\cdot\mathbf w}.
 \label{b4:CD}
\end{align}
Every exponent of $D_c$ is nonnegative, since
$\alpha\cdot\mathbf w\le4T=L_c$.

\begin{lemma}[Quartic coefficient recovery]
\label{b4:extract}
For every $\mathbf n$,
\[
 [t^{L_c}]C_{\mathbf n}(t)D_c(t)=24P_0(\mathbf n).
 \tag{B4.5}\label{b4:recovery}
\]
If $0\le n_i<b$ and $q_j$ denotes any coefficient of
$C_{\mathbf n}(t)D_c(t)$, then
\[
 |q_j|\le24H\left(1+\sum_i n_i\right)^4
       \le24H\bigl(1+31(b-1)\bigr)^4.
 \tag{B4.6}\label{b4:qbound}
\]
The product has degree at most $8T$.
\end{lemma}

\begin{proof}
The multinomial term indexed by $\gamma$ in $C_{\mathbf n}$ has
coefficient
\[
 \frac{4!}{(4-|\gamma|)!\prod_i\gamma_i!}\,\mathbf n^\gamma
\]
and exponent $\gamma\cdot\mathbf w$.  Its product with the
$\alpha$-term of $D_c$ reaches exponent $L_c$ exactly when
$\gamma\cdot\mathbf w=\alpha\cdot\mathbf w$.  Lemma~\ref{b4:code}
forces $\gamma=\alpha$, and the factorials cancel.  Summing proves
\eqref{b4:recovery}.

Distinct indices in $D_c$ have distinct exponents and coefficients
of magnitude at most $24H$.  The coefficients of $C_{\mathbf n}$
are nonnegative and their sum is $C_{\mathbf n}(1)$.  Bounding each
convolution coefficient by $24H C_{\mathbf n}(1)$ proves
\eqref{b4:qbound}.  Finally,
$\deg C_{\mathbf n}\le4T$ and $\deg D_c\le L_c=4T$.
\end{proof}

\subsection{The general mechanism behind the improvement}
\label{b4:general}

The coefficient step is not special to the particular quartic support.
It explains why the change can be much larger than a local circuit
simplification.

\begin{proposition}[Polynomial span at fixed degree]
\label{b4:generalspan}
Let $d\ge2$ and let $P\in\Z[x_1,\ldots,x_n]$ have total degree at
most $d$.  For every prime power $q\ge n+1$ there are fixed
integer weights
\[
 0=w_0<w_1<\cdots<w_n\le T\le q^d-2
\]
such that, on putting
\[
 c_\alpha=(d-|\alpha|)!\prod_i\alpha_i!\,p_\alpha,
 \qquad
 D_P(t)=\sum_{|\alpha|\le d}c_\alpha t^{dT-\alpha\cdot\mathbf w},
\]
one has the exact identity
\[
 [t^{dT}]
 \left(1+\sum_{i=1}^n x_it^{w_i}\right)^dD_P(t)=d!\,P(\mathbf x).
\]
In particular, choosing the least power of two $q\ge n+1$ gives
$T<[2(n+1)]^d$.
\end{proposition}

\begin{proof}
The Bose--Chowla construction in $\F_{q^d}$ gives a set of $q$
residue classes modulo $q^d-1$ with distinct unordered $d$-sums.
Indeed a primitive element has degree $d$ over $\F_q$;
equality of products of $d$ factors $\theta+c$ yields an
identically zero polynomial of degree at most $d-1$, exactly
as in Lemma~\ref{b4:code}.  Translate one class to zero, retain
it and any $n$ other classes, and take their representatives in
$[0,q^d-2]$.  Padding with zero gives the needed injectivity
for all multi-indices of total degree at most $d$.
The multinomial proof of Lemma~\ref{b4:extract}, with $d!$
in place of $24$, proves the displayed identity.  Finally
$q<2(n+1)$ for the stated choice of $q$.
\end{proof}

Thus at fixed $d$ the required exponent span is polynomial in
the number of source variables, whereas the standard positional
choice has span $(d+1)^n$.  This statement concerns the
coefficient-coding stage.  A full Diophantine reduction still
requires its bounds, masks, and both implications to be checked.
The following proof performs that work for the delivered quartic
instance.  The finite-field theorem is classical; no claim of
literature priority for this general use of additive codes is made.

\subsection{The new two-witness source interface}
\label{b4:interface}

For integer variables $f,g$, define the following expressions in the
displayed order:
\begin{align}
 b&=1+3f,&
 \mathcal B&=\beta_Sb^4,&
 \ell&=\sum_{i=1}^{31}\mathcal B^{w_i},
 \label{b4:source1}\\
 D_c&=D_c(\mathcal B),&
 N_0&=\mathcal B^{T+1},&
 N_1&=4\mathcal B^{8T+1},
 \label{b4:source2}\\
 N&=N_0N_1,&
 \mathscr C&=b\mathcal B^T,&
 M_c&=N_0-1-(b-1)\ell,
 \label{b4:source3}\\
 T_c&=M_c+(\mathcal B-2)\mathcal B^{L_c}N_0,
 \label{b4:source4}\\
 J_c&=2(1+g)^4D_c+\sum_{j=0}^{8T}\mathcal B^{j+1},
 \label{b4:source5}\\
 S_c&=g+J_cN_0,&
 R_c&=(S_c+T_c+1)N+T_c+1,
 \label{b4:source6}\\
 X&=(N-1)R_c,&
 Y&=N^2.
 \label{b4:source7}
\end{align}
All exponents and sum limits are fixed integers.  These expressions
are ordinary polynomials in $f,g$ with fixed integer coefficients.

\begin{theorem}[Complete source interface]
\label{b4:sourceinterface}
The following implications hold.
\begin{enumerate}
\item If $P_0$ has a natural zero, then for every positive bound $Z$
there exist $f\ge Z$ and $b\le g<\mathscr C$ such that $b$ is a
square power of two and
\[
 Y\mid\binom{2X}{X}.
 \tag{B4.7}\label{b4:binomial}
\]
\item If $f\ne0$, $0\le g<2\mathscr C$, $b$ is a square power of two,
and \eqref{b4:binomial} holds, then $P_0$ has a natural zero.
\item Already under $f\ne0$, $b$ square, and
$0\le g<2\mathscr C$, without a power-of-two assumption on $b$,
\[
 3\mid X,\qquad X>3b,\qquad Y>\max\{b,256\}.
 \tag{B4.8}\label{b4:bridgebounds}
\]
\end{enumerate}
Moreover $\mathscr C=2^{4269882T}b^{4T+1}$, with an even positive
exponent of two, so the bound has exactly the form required by
Sun's source interface.
\end{theorem}

\begin{proof}
We first establish signs and bounds before assuming a source zero or
binomial divisibility.  Since $b=1+3f$ is a square and $f\ne0$,
one has $f>0$ and $b\ge4$.  The coefficient bound
\eqref{b4:bound} implies
\begin{align}
 \frac{\mathcal B}{2}
 &>24H\bigl(1+31(b-1)\bigr)^4,
 \label{b4:firstbound}\\
 \mathcal B
 &>2+2\cdot24H(32b)^4.
 \label{b4:secondbound}
\end{align}
For the second inequality, replacing $33$ by $32$ in the bound for
$\beta_S$ leaves more than two units of margin after multiplication
by $b^4$, because $H\ge1$ and $b\ge4$.

The constant coefficient $P_0(\mathbf0)$ is positive.  Thus $D_c>0$:
its leading term is $24P_0(\mathbf0)\mathcal B^{L_c}$ and the
absolute sum of all lower terms is at most
\[
 24H\sum_{j=0}^{L_c-1}\mathcal B^j<\mathcal B^{L_c}.
\]
We also have
$D_c\le24H\sum_{j=0}^{L_c}\mathcal B^j$.
The base-$\mathcal B$ digits of $M_c$ are $\mathcal B-b$ at
positions $w_i$ and $\mathcal B-1$ at the other positions from
$0$ through $T$.  Therefore
\[
 0\le M_c<N_0,\qquad
 0<(\mathcal B-2)\mathcal B^{L_c}<N_1,\qquad
 0\le T_c<N.
 \tag{B4.9}\label{b4:blocksT}
\]

For $0\le g<2\mathscr C$, the inequalities
$2b\le\mathcal B$ and $g\in\mathbb Z$ give
$g<N_0$ and $1+g\le2b\mathcal B^T\le32b\mathcal B^T$.
Using \eqref{b4:secondbound} and the bound for $D_c$,
\begin{align*}
 0<J_c
 &<
 \frac{(\mathcal B-2)+\mathcal B}{\mathcal B-1}
       \mathcal B^{8T+1}\\
 &=2\mathcal B^{8T+1}<N_1.
\end{align*}
Consequently $0\le S_c<N$.  These are unconditional bounds on
the stated source range.

Because $\beta_S$ is an even power of two and $b\equiv1\pmod3$,
we have $\mathcal B\equiv1\pmod3$ and
$N=4\mathcal B^{9T+2}\equiv1\pmod3$.  Hence $3\mid X$.
Also $R_c\ge N+1$, $N>\mathcal B\ge b$, and $N\ge4$, so
$X=(N-1)R_c>3b$ and $Y=N^2>\max\{b,256\}$.
This proves \eqref{b4:bridgebounds}.

Now assume additionally that $b$ is a power of two.  Then
$\mathcal B,N_0,N_1,N$ are powers of two.  For nonnegative integers
$a,d$, write $\tau_2(a,d)=0$ when their binary addition has no
carries, equivalently when their bitwise AND is zero.  The mask
description gives
\begin{multline}
 0\le g<2\mathscr C,\quad \tau_2(g,M_c)=0\\
 \Longleftrightarrow\quad
 g=\sum_{i=1}^{31}n_i\mathcal B^{w_i}
 \quad\hbox{for some }0\le n_i<b.
 \tag{B4.10}\label{b4:mask}
\end{multline}
Indeed the digit $\mathcal B-b$ has ones exactly in the high binary
positions of one $\mathcal B$-digit, whereas $\mathcal B-1$ has
ones in every position.  The only free bits of $g$ are therefore
the low $\log_2b$ bits at the distinct slots $w_i$.  The bound
$g<N_0$ excludes all higher slots.  Conversely every displayed
code has the required bit pattern and is less than $\mathscr C$.
No superincreasing property of the weights is used.

For decoded $g$, write
\[
 (1+g)^4D_c
 =\sum_{j=0}^{8T}q_j\mathcal B^j
\]
using the formal coefficients of Lemma~\ref{b4:extract}, and put
\[
 J'=\sum_{j=0}^{8T}
          \left(q_j+\frac{\mathcal B}{2}\right)\mathcal B^j.
\]
Equations \eqref{b4:qbound} and \eqref{b4:firstbound} show that
every displayed digit of $J'$ is an integer strictly between
$0$ and $\mathcal B$.  In particular this is its actual
base-$\mathcal B$ expansion, even when some $q_j$ are negative.
Since $J_c=2J'$,
\[
 \tau_2(J_c,(\mathcal B-2)\mathcal B^{L_c})=0
 \ \Longleftrightarrow\
 \tau_2\!\left(J',
       \left(\frac{\mathcal B}{2}-1\right)\mathcal B^{L_c}\right)=0.
\]
The mask on the right occupies precisely the low bits of the
digit at $L_c$.  A digit in $[1,\mathcal B)$ avoids it exactly
when it equals $\mathcal B/2$.  By \eqref{b4:recovery}, this is
equivalent to $q_{L_c}=24P_0(\mathbf n)=0$.

The two block decompositions, together with $S_c,T_c<N$, give
\[
 \tau_2(S_c,T_c)=0
 \ \Longleftrightarrow\
 \tau_2(g,M_c)=0
 \ \land\
 \tau_2(J_c,(\mathcal B-2)\mathcal B^{L_c})=0.
 \tag{B4.11}\label{b4:blocksplit}
\]
Sun's binary two-block lemma \cite[Lemma~2.2]{Sun2021} states,
for $0\le S_c,T_c<N$ with $N$ a power of two and the exact
$R_c,X$ above, that
\[
 \tau_2(S_c,T_c)=0
 \quad\Longleftrightarrow\quad
 N^2\mid\binom{2X}{X}.
 \tag{B4.12}\label{b4:Kummer}
\]
Its hypotheses have all been established before invoking it.
Combining \eqref{b4:mask}--\eqref{b4:Kummer} proves soundness.

For completeness, suppose $P_0(\mathbf n)=0$.  Choose $s$ so large
that $b=2^{2s}>\max_i n_i$ and $(b-1)/3\ge Z$.  Since
$4^s\equiv1\pmod3$, $f=(b-1)/3$ is an integer.  The source
zero is not the all-zero vector, so
$g=\sum_i n_i\mathcal B^{w_i}\ge\mathcal B\ge b$.
The usual geometric-sum estimate gives $g<b\mathcal B^T=\mathscr C$.
The mask and coefficient identities yield
\eqref{b4:binomial}.  Finally the asserted monomial formula for
$\mathscr C$ follows directly from $\mathcal B=\beta_Sb^4$.
\end{proof}

\subsection{The complete outer polynomial}
\label{b4:outer}

We now give the full outer formula, so that the artifact is specified
by fixed constants, finite sums, and ordinary polynomial operations.
Only the ten lowercase coordinates listed in
Theorem~\ref{b4:main} are unknowns.  Every other symbol is a circuit
register.

Starting with the new $b,X,Y,\mathscr C$ above, set
\begin{align}
 L_s&=lY,& U&=2L_sX,\notag\\
 V&=4gwY,& W_s&=bw,
 \label{b4:bridge1}\\
 A&=U(V+1),& B_s&=2X+1,\notag\\
 C&=B_s+(A-2)h,& P_n&=U^2V,
 \label{b4:bridge2}\\
 K&=X+1+k(P_n-2),&
 S_s&=2A-5,
 \label{b4:bridge3}\\
 T_s&=3W_sC-2(W_s^2-1),&
 R_2&=(P_n^2-4)K^2+4,
 \label{b4:bridge4}
\end{align}
\begin{equation}
 \Delta=8\mathscr C^3gK^2
       -g^2\{32(C-KL_s)^2\mathscr C^3+g^2K^2\},
 \label{b4:bridge5}
\end{equation}
\begin{align}
 O&=f^2l^2\Delta,&
 \xi&=zx.
 \label{b4:bridge6}
\end{align}
The modified Lucas certificate is
\begin{align}
 D_s&=(A^2-4)C^2+4,&
 E&=C^2D_s\xi,
 \label{b4:lucas1}\\
 F_s&=4(A^2-4)E^2+1,&
 G_*&=2(A+2)(A-2)^2E^2-1,
 \label{b4:lucas2}\\
 \Omega&=16E,&
 H_*&=C+(B_s-C)F_s+\Omega F_sy,
 \label{b4:lucas3}\\
 I_*&=(G_*^2-1)H_*^2+1.
 \label{b4:lucas4}
\end{align}
These are the same bridge and Lucas formulas as in package~31
\cite{Package31}, using Sun's modified certificate and interpolation
theorem \cite{Sun2026}.

For the divisibility and nonvanishing merger, define
\begin{align}
 r&=T_s+S_sy,&
 \Sigma&=\xi S_s^5,
 \label{b4:merge1}\\
 \Theta&=(r^3-rS_s^2-S_s^3)(r^2+23S_s^2).
 \label{b4:merge2}
\end{align}
The three radicands and their separating weights are
\begin{align}
 r_0&=b\{1+(3O-4)b^2z^2\},&
 r_1&=R_2,&
 r_2&=D_sF_sI_*,
 \label{b4:radicands}\\
 \omega_0&=1,&
 \omega_1&=r_0,&
 \omega_2&=r_0(P_n^2+K^2+3).
 \label{b4:weights}
\end{align}
Finally put
\begin{align}
 Z_c&=\Sigma m-\Theta,&
 A_i&=\Sigma^2\omega_i^2r_i\quad(0\le i\le2),
 \label{b4:norm1}\\
 V_n&=Z_c^2+A_2-A_0-A_1,
 \label{b4:norm2}\\
 W_n&=V_n^2+4Z_c^2A_2-4A_0A_1,
 \label{b4:norm3}\\
 \boxed{F_{B_4}}&\boxed{{}=W_n^2-16V_n^2Z_c^2A_2.}
 \label{b4:norm4}
\end{align}
The requested equation is $F_{B_4}=0$.  No radical in the proof below
is part of its definition.  There are no divisions or variable
exponents in \eqref{b4:source1}--\eqref{b4:norm4}.

\subsection{Why the finalizer retains both implications}
\label{b4:finalizer}

For clarity we record the inherited finalizer arguments, with their
attribution to package~31 \cite{Package31}.  They ensure that the new
source interface is being composed with the precise outer hypotheses.

\begin{lemma}[The quintic merger]
\label{b4:quintic}
The polynomial
$\pi(v)=(v^3-v-1)(v^2+23)$ has a root modulo every positive integer
and no rational root.  If $S_s\ne0$ and $\xi\ne0$, then
$\Sigma\mid\Theta$ implies $S_s\mid T_s$.
Conversely, when $S_s\mid T_s$, a residue class of $y$ modulo
$|\xi|$ makes $\Sigma\mid\Theta$.  For every integer assignment
with $S_s\ne0$, $\Theta\ne0$.
\end{lemma}

\begin{proof}
The cubic has discriminant $-23$.  At $p=23$ it has the simple
root $3$.  At an odd prime $p\ne23$, either it has a simple root
or it is irreducible.  In the latter case Frobenius is an even
three-cycle, so the discriminant is a square modulo $p$ and
$v^2+23$ has a simple root.  Hensel lifting handles all odd prime
powers.  Since $-23\equiv1\pmod8$, the quadratic has a root in
the $2$-adic integers as well.  The Chinese remainder theorem
gives a root modulo every positive integer.

The rational-root theorem excludes rational roots of the cubic,
whose only possible integer roots are $\pm1$; the quadratic is
positive on $\mathbb Q$.  Consequently
$\Theta=S_s^5\pi((T_s+S_sy)/S_s)$ cannot vanish.

If $S_s\mid T_s$, write $T_s=S_sq$ and choose
$y\equiv c-q\pmod{|\xi|}$ for a root $c$ of $\pi$ modulo $|\xi|$.
Then $\Theta=S_s^5\pi(q+y)$ is divisible by $\Sigma$.
Conversely suppose that a prime satisfies
$a=v_p(S_s)>d=v_p(T_s)$.  Then $v_p(r)=d$; the leading term of
the homogeneous cubic factor has uniquely smallest valuation
$3d$, and that of the quadratic factor has uniquely smallest
valuation $2d$.  Thus $v_p(\Theta)=5d<5a$, contradicting
$S_s^5\mid\Theta$.  Therefore $S_s\mid T_s$.
\end{proof}

\begin{lemma}[Positivity and separating weights]
\label{b4:separation}
For the displayed polynomial expressions, if $z\ne0$ then
$|r_0|>1$ and
\[
 \sqrt{|r_1|}\le P_n^2+K^2+2.
 \tag{B4.13}\label{b4:r1bound}
\]
Any rational signed sum
$\pm\sqrt{r_0}\pm\omega_1\sqrt{r_1}\pm\omega_2\sqrt{r_2}$
forces every $r_i$ to be an integer square.
Moreover, $r_0$ square and $z\ne0$ imply that $b$ is a positive
square and $O>0$.
\end{lemma}

\begin{proof}
The congruence $b\equiv1\pmod3$ excludes $b=0,-1$.  If
$r_0=\pm1$, its factor $b$ must be $1$, hence $f=O=0$ and
$r_0=1-4z^2$, a contradiction.  If $r_0=0$, its bracket vanishes,
forcing $|b|=|z|=1$ and $3O-4=-1$; again $b=1$ gives $O=0$,
a contradiction.  This proves $|r_0|>1$.

For \eqref{b4:r1bound}, use
$|r_1|\le(P_n^2+4)K^2+4$.  Subtracting this bound from
$(P_n^2+K^2+2)^2$ leaves
$P_n^4+K^4+P_n^2K^2+4P_n^2\ge0$.
It follows that
\[
 |\omega_1|>\sqrt{|r_0|},\qquad
 |\omega_2|>\sqrt{|r_0|}
                    +|\omega_1|\sqrt{|r_1|}.
 \tag{B4.14}\label{b4:strictweights}
\]
In the multiquadratic splitting field, any automorphism changing a
nonzero term of a rational signed sum has a largest changed index.
The magnitude of that term is at least the magnitude of its weight,
since its nonzero radicand is an integer.  By
\eqref{b4:strictweights}, all earlier terms together are too small
to cancel its change.  This contradiction shows that every root
term is rational.  A rational square root of an integer is integral.
The argument also covers negative candidate radicands by complex
absolute values, and zero radicands require no recovery.

Finally the two factors of $r_0$ are coprime, since the bracket
is $1$ modulo $b$.  If their product is a nonzero square, each
factor is a signed square with the same sign.  A negative square
cannot be $1$ modulo $3$, so $b$ is a positive square and its
bracket is a square.  If $O\le0$ and $bz\ne0$, that bracket
would be at most $1-4=-3$.  Hence $O>0$.
\end{proof}

\begin{proof}[Proof of Theorem~\ref{b4:main}: equivalence]
The polynomial identity underlying the finalizer is
\[
 F_{B_4}
 =\prod_{\varepsilon\in\{-1,1\}^3}
   \left(Z_c+\sum_{i=0}^2\varepsilon_i\sqrt{A_i}\right).
 \tag{B4.15}\label{b4:normidentity}
\]
It follows by successively multiplying conjugate factors, or by
expanding the $35$ terms of the three-radicand norm.  The delivered
checker proves this exact formal identity independently.

Suppose $F_{B_4}=0$.  The integer $S_s=2A-5$ is odd and nonzero,
so Lemma~\ref{b4:quintic} gives $\Theta\ne0$.  If $\Sigma=0$,
then every $A_i=0$ and \eqref{b4:norm4} equals $\Theta^8\ne0$.
Thus $\Sigma\ne0$, and $\xi=zx\ne0$.  Divide a vanishing factor
of \eqref{b4:normidentity} by $\Sigma$.  Lemma~\ref{b4:separation}
makes every $r_i$ a square; the resulting signed sum of roots is
an integer, so $\Sigma\mid\Theta$.
Lemma~\ref{b4:quintic} gives $S_s\mid T_s$, and
Lemma~\ref{b4:separation} gives $b$ square and $O>0$.
These, $\xi\ne0$, $R_2$ square, and $D_sF_sI_*$ square are
exactly the hypotheses of Sun's reverse bridge
\cite[Proposition~6.2]{Sun2026}.  It yields the source range,
the power-of-two condition, and binomial divisibility required
by Theorem~\ref{b4:sourceinterface}.  Hence $P_0$ has a natural
zero, and \eqref{b4:endpoint} yields a finite ZFC refutation.

Conversely suppose $P_0$ has a natural zero.
Theorem~\ref{b4:sourceinterface} supplies arbitrarily large $f$;
choose $f\ge2$.  Sun's positive bridge
\cite[Proposition~6.1]{Sun2026} supplies the other initial
witnesses with $b$ square, $R_2$ square, $S_s\mid T_s$,
$C=u_{B_s}(A)$, and $O>0$.  As
$O=f^2l^2\Delta$ is a positive integer, $O\ge4$.
The positive nonsquare integer $d=3O-4\equiv2\pmod3$
has a Pell solution with its second coordinate a positive
multiple of $b$: take a positive Pell unit to a power equal
to the order of its class modulo $b$.  Write that second
coordinate as $bz$ with $z>0$.  Then $r_0$ is a square.

Sun's full residue-class interpolation theorem
\cite[Theorem~4.1]{Sun2026} supplies a positive $\xi$ divisible
by this $z$, and allows $y$ in the residue class required by
Lemma~\ref{b4:quintic}, while making $D_sF_sI_*$ a square.
Set $x=\xi/z$.  Thus all $r_i$ are squares,
$\Sigma\ne0$, and $\Sigma\mid\Theta$.  Choose the integer $m$
so that a selected signed factor of \eqref{b4:normidentity}
vanishes.  This gives $F_{B_4}=0$.

All new steps have been proved above.  The substantial imported
machine/compiler and Lucas results are explicitly identified;
their ordinary arithmetic proofs, not the soundness of any
proof assistant, supply the asserted ZFC-provability.
\end{proof}

\subsection{Circuit identity, degree, and reproducibility}
\label{b4:verification}

The sparse coefficient evaluator uses the division-free identity
\[
 \mathcal B^{4T-\alpha\cdot\mathbf w}
 =(\mathcal B^T)^{4-|\alpha|}
       \prod_{i=1}^{31}
          (\mathcal B^{T-w_i})^{\alpha_i}.
 \tag{B4.16}\label{b4:horner}
\]
All exponents are nonnegative.  A multivariate Horner scheme shares
the powers and coefficient recipes.  Equation~\eqref{b4:horner}
is valid for arbitrary integer assignments, including
$\mathcal B=0$; no division by $\mathcal B$ is concealed.
The geometric sum in \eqref{b4:source5} uses binary identities
for fixed powers and fixed-length sums.

The main circuit contains $485$ additions, $25$ subtractions, and
$1694$ multiplications, for $2204$ operations.  There are $79$
fixed coefficient leaves after reuse and literal simplification.
The general source-degree bounds in this instance are
\[
 \deg X=108T+24,\qquad
 \deg Y=72T+16,\qquad
 \deg\mathscr C=4T+1.
 \tag{B4.17}\label{b4:sourcedegrees}
\]
Propagating degree bounds through the circuit yields
\[
 \deg F_{B_4}\le113856T+27064=107286763576.
 \tag{B4.18}\label{b4:degreeupper}
\]
To certify attainment rather than just an upper bound, substitute
one indeterminate for all ten unknowns.  With the actual fixed
coefficient payloads, the coefficient at the bound in
\eqref{b4:degreeupper} is $16$ modulo $101$.  It is also $599$
modulo $1009$ and $622356$ modulo $1000003$.  A single nonzero
such coefficient proves attainment, so the stated total degree
is exact.  The checker performs exact leading-coefficient
arithmetic modulo the primes; this is not inference from
random numerical agreement.

For comparison, package~31's positional construction has exact degree
\[
 126496\cdot5^{31}+27064
 =589042901992797851562527064
\]
and $1442$ operations.  The new degree is smaller by a factor of
approximately $5.49\times10^{15}$, at the cost of $762$ additional
arithmetic operations.  Thus the degree improvement is substantive,
while the operation-count tradeoff remains explicit.

The self-contained commands are
\begin{quote}
\code{python3 b4_code.py}\\
\code{python3 build_b4_sun10.py}\\
\code{python3 verify_b4_sun10.py}
\end{quote}
from \code{ten_integer/}.  They use only the Python standard
library.  In addition to the degree proof, they regenerate the
finite field and all $52360$ code sums, reconstruct all $402$
weighted coefficient records, check the fixed coefficient payloads
and bounds, and expand the norm identity exactly.  A separate
scalar transcription evaluates the literal $402$-term coefficient
sum and all named registers modulo three primes, independently
of the Horner circuit.  Those modular transcription checks are
regression evidence; the preceding lemmas supply the universal
equivalence.

The archive also includes four finalizer variants using more generic
weights or the literal Sun four-square merger.  They share the same
source adapter and fixed tuple.  A clean isolated rebuild reproduces
every generated JSON file byte for byte.  Source pins, coefficient
bindings, licenses, and inherited attribution notices accompany
the construction.

\section{Nonzero label codes and smaller fixed coefficients}
\label{sec:coefficient-refinement}

The variable-compression construction and the coefficient construction
address different measures.  This section gives a new coefficient tuple
and binds two complete, moderate-degree equations to it.  It does not
silently substitute these coefficients into a low-variable form whose
index constants were computed from another parent polynomial.

\subsection{The new tuple and its measured improvement}

The baseline is package 36 of the reviewed repository commit.  It uses
943 Wang tiles, a radix $\rho=2^{115}$, cell length $8801$, and
$B=2^{1012115}$.  The new construction uses the identical tile language,
but has
\begin{equation}
 b=107,\qquad L=6385,\qquad d=bL=683195,\qquad
 \rho=2^{107},\qquad B=\rho^L=2^{683195}.
 \label{eq:nonzero-base}
\end{equation}
All four resulting coefficients are materialized, with their exact
unsigned little-endian payloads and hashes included in the archive.
The comparison is:
\begin{center}
\begin{tabular}{lrr}
\hline
Measure & Package 36 & New tuple\\
\hline
Base exponent $\log_2 B$ & $1{,}012{,}115$ & $683{,}195$\\
Largest coefficient, bits & $1{,}016{,}946$ & $688{,}546$\\
Four raw payloads, bytes & $506{,}664$ & $342{,}269$\\
\hline
\end{tabular}
\end{center}
These decrease by $32.50\%$, $32.29\%$, and $32.45\%$, respectively.
There are two mathematical changes: an occupancy bit on each edge word
is replaced by nonzeroness of that whole word; and exact support
isolation permits a much shorter synchronization layout.

\subsection{Occupancy requires a nonzero word, not an extra bit}

\paragraph{Nonzero-code lemma.}
\emph{Suppose a local selector/copy clause permits either an empty cell,
whose edge words are all zero, or one selected tile with its four encoded
edge words.  Suppose the horizontal and vertical label encodings are
injective and never zero, and corresponding words agree at neighboring
cells.  A Start tile in a cyclic word whose spatial shift is one forces
every cell to be occupied and decodes the original tile matching rules.}

Indeed, a neighbor of an occupied cell has a matching nonzero word and
therefore cannot be empty.  Iteration of spatial step one visits the
whole cyclic word.  Once occupancy is established, injectivity decodes
the original labels.  Conversely, any tile matching supplies the encoded
words and satisfies all bit comparisons.  This argument does not require
a bit position shared by every nonzero word.

The inherited horizontal colors are $0,\ldots,249$, and the vertical
colors are $0,\ldots,438$.  Encode every label $x$ as $x+1$.  Horizontal
words now use eight bits and vertical words nine.  This is optimal
\emph{within the fixed-length, separate-edge encoding interface with zero
reserved for emptiness}: $m$ labels require $2^\ell-1\ge m$.  The bounds
are therefore $\ell_H\ge8$ and $\ell_V\ge9$, attained here.  This is not
a lower bound on general Diophantine constructions.

\subsection{The entire finite recipe}

Keep the ordered 943 Wang tiles, with Start at selector $0$ and Forbidden
at selector $1$.  The language and literal 430-action K215 table are
unchanged.  The separately supplied tile list has SHA-256
\begin{center}\small\ttfamily
66b55f1d11fca9fefc93b12fcfc8b4cabb75\\
b08e9fd03311f9e894a1c5e5e81f.
\end{center}
Use the following slot data:
\begin{center}
\begin{tabular}{lrrr}
\hline
Slot $s$ & Offset $o_s$ & Width $\ell_s$ & Field $p_s$\\
\hline
West & 1023 & 8 & 0\\
North & 1013 & 9 & 8\\
East & 1003 & 8 & 17\\
South & 983 & 9 & 25\\
\hline
\end{tabular}
\end{center}
The native support $E$ is the disjoint union of $\{0,\ldots,942\}$,
the 34 copy positions $o_s+j$, the dummies $1063,1064$, and the anchors
$1022,1042,2107$.  Hence $|E|=982$.  Set $w_j=2^{10+2j}$ for
$0\le j<37$.  For tile $i$, with label $x_{i,s}$ in slot $s$, define
\[
 c_i=1+\sum_{s=0}^3\sum_{j=0}^{\ell_s-1}
        \operatorname{bit}_j(x_{i,s}+1)w_{p_s+j}
        +w_{34}+w_{35}+w_{36}.
\]
At the copy positions set $c_{o_s+j}=w_{p_s+j}$; at the three anchors
use $w_{34},w_{35},w_{36}$, in the displayed order; and set both dummy
coefficients to zero.  Define the following fixed integers:
\begin{align*}
 \mu&=1022+\sum_{j=0}^{36}w_j,\\
 Z&=(2^{105}-1)\sum_{e=6376}^{6383}\rho^e
                +(2^{30}-1)\rho^{6384},\\
 D_C&=\rho^{30}+\rho^{70}+\rho^{20}+\rho^{1085}
       +\sum_{e\in E}c_e(\rho^{3202-e}+\rho^{4268-e}),\\
 D_R&=\rho^{50},\\
 M_C&=B-1-\sum_{\substack{e\in E\\e\ne1}}\rho^e-2\rho^{1063},\\
 M_{F,\mathrm{native}}&=4+\mu(\rho^{3202}+\rho^{4268})
       +\sum_{j=0}^8\rho^{1013+j}+\sum_{j=0}^7\rho^{1073+j}
       +\rho^{1042}+\rho^{2107}+Z.
\end{align*}
The four actual ports are
\begin{equation}
 \boxed{\beta=B-1,\qquad K=D_C+BD_R,\qquad
 M_C,\qquad M_F=M_{F,\mathrm{native}}+B-1.}
 \label{eq:nonzero-ports}
\end{equation}
In particular $M_F$ already includes its shift.  All sums and powers in
this recipe are fixed finite arithmetic; none is an existential variable
or a computation that waits for a refutation.

\subsection{Why the new finite interface is sound and complete}

Let $n$ be the selector count in one native cell.  Before inspecting any
payload field, $0\le n\le943$, and the lowest ten bits of the clause sum
are precisely $n$.  Disjointness with $1022$ forces $n=0$ or $n=1$.
Only then does each radix-four payload field become a sum in
$\{0,1,2\}$.  Testing its low bit forces equality of the selected tile
bit and its copy, and forces each anchor to equal $n$.  Thus payload
carry-freeness is not assumed before selector extraction.

The complete comparison identities are:
\begin{center}\small
\begin{tabular}{llll}
\hline
Target & Current source via $D_C$ & Spatial source & Temporal source\\
\hline
$1013+j$, $0\le j<9$ & $983+j$ & none & $1013+j$\\
$1073+j$, $0\le j<8$ & $1003+j$ & $1023+j$ & none\\
$1042$ & $1022$ & none & $1042$\\
$2107$ & $1022$ & none & $2107$\\
\hline
\end{tabular}
\end{center}
At each clause band $T=3202,4268$, the coefficient of $\rho^{T-e}$ in
$D_C$ is exactly $c_e$ for every $e\in E$, and $e+50\ne T$.  The
generator and independently written auditor check all
$19\cdot982\cdot3+2\cdot982\cdot2=59{,}902$ scalar identities.  Their
linearity covers every native assignment, including the allowed dummy
digits.  No additional bare or reflected-clause contribution is omitted.

Write $S=E\setminus\{2107\}$.  Its maximum is $1064$.
The difference $1066$ is absent from $E-E$, since pairs in $S$ have
smaller span and $2107-1066=1041\notin E$.  Because $L>2\max E$,
\[
 ((3202-E)\bmod L)\cap((4268-E)\bmod L)=\varnothing.
\]
At least one clause band is clean for every binary rotation.  Decode
the distinguished Start cell first: its selector is already one, so its
source anchor at $1022$ is one and both tested temporal anchors must be
one.  This step does not presuppose global occupancy.

The source's parity argument excludes nonzero binary remainders: for
remainders $1,\ldots,b-2$ every rotated native digit is even, and for
$b-1$ only the single upper-dummy spill can be odd.  Two tested anchors
cannot both be odd in either case.  The remaining whole-radix alignment
is unique, since the difference $2107-1042=1065$ occurs exactly once in
$E-E$.  Thus
\[
 ((1042-E)\bmod L)\cap((2107-E)\bmod L)=\{0\}.
\]
After alignment, the parity comparisons give equality of corresponding
bits.  The nonzero-code lemma then propagates occupancy and recovers
the unchanged marked Wang tiling.  Conversely every intended tiling
supplies these words and anchors, so the change loses no tableau.

The inherited conservative mass bound is
\[
 \mathcal M=(982+2)\left(2\sum_{e\in E}c_e+6\right).
\]
Its bit length is 104, while $4\mathcal M+8$ has bit length 106;
therefore $4\mathcal M+8<2^{107}$.  Also $2\mu+4<2^{107}$ and
$2(2^{105}-1)+4<2^{107}$.  These are the inherited uniform carry and
strict-slack bounds.  The maximum unshifted current exponent is
$4268+2107=6375$, and the spatial maximum is $50+2107=2157$.
All padding positions $6376,\ldots,6384$ are consequently zero on every
intended aligned field, including safe dummy variations.

The unpadded high-mask population is $2\cdot46+9+8+2+1=112$.
Padding contributes $8\cdot105+30=870$, giving
\[
 \operatorname{pc}(M_{F,\mathrm{native}})=982,\qquad
 \operatorname{pc}(M_C)=d-982.
\]
Moreover $M_C\equiv2\pmod4$, $M_{F,\mathrm{native}}\equiv4\pmod8$,
and both native masks lie strictly between zero and $B-1$.  The origin
test is unchanged because every exponent in $D_C,D_R$ is positive.

\subsection{Controlling the actual Pell rank}

Changing $d$ requires a new actual-rank controller.  We have
\[
 d=5\cdot107\cdot1277,
 \qquad \lambda=135257\text{ prime},\qquad k=\lambda-1=135256.
\]
The orders of $B$ modulo $5,107,1277,\lambda$ are respectively
$4,106,1276,67628$.  Hence
$\operatorname{ord}_d(B)=67628\mid k$, while
\[
 B^k\bmod\lambda^2=13998558473\ne1,
 \qquad v_\lambda(B^k-1)=1.
\]
For $H_0=1+BD_C+B^2(D_R+1)$, the exact unit certificate is
\begin{center}
\begin{tabular}{rrrr}
\hline
$p$ & $B\bmod p$ & $(B^2-1)\bmod p$ & $H_0\bmod p$\\
\hline
5 & 3 & 3 & 2\\
107 & 88 & 39 & 58\\
1277 & 942 & 1125 & 956\\
135257 & 55321 & 88158 & 66933\\
\hline
\end{tabular}
\end{center}
Choose independently padded width $w_0$ and height $h_0$ to be
sufficiently large powers of $\lambda$, put $N=w_0h_0=\lambda^n$,
$q=B^N$, and require $N>d\lambda^2$ and $h_0>\lambda$.
Then $q\equiv B^{w_0}\equiv B\pmod{d\lambda}$.  Flipping a low
dummy at cell $j$ changes the actual packed rank by
\[
 \delta R=-\Gamma B^j,\qquad
 \Gamma=\rho^{1063}(q^2-1)
             [1+q(D_C+BD_R+B^{w_0})].
\]
The table proves $\gcd(\Gamma,dN)=1$.

For completeness, put $T_0=N/\lambda$ and $C_0=d\lambda$.
Every residue modulo $dN=C_0T_0$ is a Boolean subset sum of
\[
 \{B^{2+kj}:0\le j<T_0\}\cup\{B^3\}.
\]
Indeed the first $T_0$ powers of $B^k$ are exactly the residues
$1+C_0z$.  Normalize a target by $B^{-2}$ and subtract
$\epsilon B$, $\epsilon\in\{0,1\}$, to make it nonzero modulo
$\lambda$.  Its least residue $a$ modulo $C_0$ satisfies
$1\le a<C_0<T_0$ and $\gcd(a,T_0)=1$.  Writing the target as
$a+C_0u$, solve
\[
 az_0+\frac{a(a-1)}2\equiv u\pmod{T_0}.
\]
The $a$ distinct orbit entries with additive coordinates
$z_0,z_0+1,\ldots,z_0+a-1$ have the desired sum.  Include $B^3$
if $\epsilon=1$.  The greatest selected position is at most
$N-N/\lambda-\lambda+3$, and $w_0<N/\lambda$; hence no required
spatial or temporal translate wraps or touches the origin.

Starting from rank $R_0$, choose the subset for
$(R_0-dw_0)\Gamma^{-1}$ modulo $dN$.  It gives
$R\equiv dw_0\pmod{dN}$, exactly the temporal rotation required by
$X=2^R$.  The strict positive slack survives every safe dummy choice.
The inherited Pell converse supplies the remaining witnesses.  Prime
powers are a completeness choice only; soundness does not impose them
on arbitrary solutions.

Both $d$ and $N$ are odd, and the canonical word is $1$ modulo $4$.
The balanced populations and low congruences therefore preserve the
inherited identities $\operatorname{pc}(R)=3dN+2$ and $R\equiv3\pmod4$.
In particular $X/q$ remains a positive square, and the inherited odd
Pell multipliers and odd half-binomial scale remain available.  This
also verifies the completeness promises of the compatible signed
arithmetic adapters.

\subsection{Two concrete equations bound to the new tuple}

To make the coefficient result a complete mathematical instance, the
archive supplies the old arithmetic graphs with all four ports rebound
to \eqref{eq:nonzero-ports}.  The 15 positive variables for the first are
\[
 J,F,\alpha,z,f,h,i,T,s,w,\tau,\eta,\zeta,y,\gamma.
\]
The following are polynomial abbreviations, not extra unknowns:
\begin{align*}
 q&=\beta J+1,& C&=q-F-\alpha,\\
 X&=wq,& Y&=sq^3,\\
 E_0&=XY,& k_0&=\eta+\zeta,\\
 c&=k_0Y+\eta,& a&=E_0+Y,\\
 H&=4a+3,& \Delta&=a^2+H,\\
 D&=X+ac+\gamma H,& G&=(q-1)(q-F)+\alpha.
\end{align*}
Also set
\begin{align*}
 R&=G(q^2-1)+(M_C+qM_F)J,\\
 V&=c(Tf-1)-Rf^2,\qquad K_a=(i\Delta c^2)^2.
\end{align*}
Set
\begin{align*}
 N_0&=\tau^2-(XY^2)(XY^2+1)k_0^2,\\
 N_1&=D^2-\Delta c^2,\\
 N_a&=K_a(V^2-y^2)+y^2,\\
 N_k&=k_0-hE_0-R,\\
 N_t&=(K+w)C+(q-F)-z(q-1),\\
 N_s&=\Delta f^2-K_a.
\end{align*}
Then the fixed equation is
\begin{equation}
 \boxed{P_{67}^{\mathrm{nz}}=N_0N_1N_aN_kN_tN_s-\Delta=0.}
 \label{eq:nonzero-p67}
\end{equation}
Its delivered graph has 67 binary arithmetic operations, comprising
39 multiplications and 28 additions/subtractions.  Its exact degree is
155.  The crucial cancellation is the polynomial identity
\[
 N_1=2ac(X+\gamma H)+(X+\gamma H)^2-Hc^2,
\]
which gives degree 18.  The six factor degrees are
$22,18,60,7,2,46$.  For example, the degree-155 monomial
\[
 J^{92}h\gamma i^4\eta^{13}w^{13}s^{26}T^2f^2z
\]
has coefficient $8\beta^{92}\ne0$.  Thus the degree remains exact after
the new fixed tuple is bound.

For the second presentation replace $T$ by two positive unknowns $o,j$,
so that there are sixteen unknowns in total, and retain the same outer
abbreviations.  Put
\[
 L_0=E_0k_0Y,\quad V_0=of-c,\quad U_0=jc-R,
 \quad I=ic^2,\quad A_0=\Delta(f^2-1).
\]
The seven equations are
\begin{equation}
\begin{aligned}
 \tau^2-L_0(L_0+k_0)&=1,& D^2-\Delta c^2&=1,\\
 A_0(V_0^2-y^2)+y^2&=1,& k_0-hE_0-R&=1,\\
 (K+w)C&=F+z(q-1),& I^2&=A_0,\\
 U_0&=V_0.&&
\end{aligned}
\label{eq:nonzero-system59}
\end{equation}
Their residual degrees are $22,18,24,7,2,22,6$.  The graph evaluates
the comparisons using 59 arithmetic operations, comprising 33
multiplications and 26 additions/subtractions; comparisons are free
under the repository convention.  These are inherited arithmetic
graphs, not new operation records.

\paragraph{Bound-instance theorem.}
\emph{ZFC proves that each of
\eqref{eq:nonzero-p67} and \eqref{eq:nonzero-system59} has no solution
in its stated positive-integer domain if and only if $\operatorname{Con}(\mathrm{ZFC})$.}

The inherited Pell/arithmetic proofs equate solutions of these graphs
with the blank marked native compiler problem, provided its complete
four-port interface holds.  The preceding argument establishes that
interface for the new tuple, including soundness, actual-rank tuning,
strict positive witnesses and canonical signed promises.  The unchanged
Wang-language theorem equates a marked tiling with a finite blank K215
halting run, and the K215 calibration equates such a run with a finite
ZFC refutation.  Composition and negation prove the displayed theorem.
This is a proof of an equivalence, with the inherited mathematical
dependencies made explicit; it is neither a decision of consistency nor
a new end-to-end kernel check.

\subsection{Reproducible evidence and limits}

The directory \texttt{coefficient\_backend/} contains the full recipe,
literal raw and gzip coefficients, machine actions, independently pinned
tiles, preserved notices, inherited proof sources, and the new proof.
Its builder and auditor require only the Python standard library.
The independent review found no gap in the nonzero-code bootstrap or
the new synchronization support.

The concrete circuit verifier proves that every arithmetic graph,
witness list and comparison is unchanged from its pinned source, and
that every coefficient port is fixed to the declared payload hash.
It reproduces exact degree bounds using the main-norm cancellation,
and proves degree attainment by a nonzero leading coefficient on a
specified line modulo $1000000007$.  This latter computation is an
exact lower-degree certificate, not probabilistic identity testing.
All 67 or 59 gates and all four ports are live.

The archived \texttt{PROOF.md} gives the complete adapter derivation,
including the Boolean subset-sum controller.  Finite checks verify its
explicit constants and support identities; they do not by themselves
prove the unbounded Pell, tiling or K215 calibration theorems.  The
comparison ten-integer compression retains its original parent tuple
and fixed indices.  Section~\ref{composed:section} completes the
rebinding as a distinct, fully regenerated ten-integer instance.

\section{A complete combined instance}
\label{composed:section}

The two improvements can be composed without changing the number of
unknowns or the arithmetic graph.  The separate directory
\code{ten_integer_composed/} performs this composition explicitly:
it retains the literal quartic parent and the $B_4$ code, substitutes all
four ports of the nonzero683195 backend together, recomputes every source
coefficient bound, and materializes a new $\beta_S$.
The original \code{ten_integer/} instance remains available for the
isolated degree comparison.

\begin{corollary}[The combined ten-integer equation]
\label{composed:main}
The fixed circuit
\code{ten_integer_composed/generated/b4_signed10.json},
with its local \code{data/INSTANCE.json} binding, defines a polynomial
$F^{\mathrm{nz}}_{B_4}\in\Z[f,g,h,k,l,w,x,y,z,m]$ satisfying
\[
 \mathrm{ZFC}\ \vdash\
 \left[
   \bigl(\exists\mathbf u\in\Z^{10}:
            F^{\mathrm{nz}}_{B_4}(\mathbf u)=0\bigr)
   \longleftrightarrow\operatorname{Ref}_{\mathrm{ZFC}}
 \right].
\]
Thus its lack of integer zeros is equivalent in ZFC to
$\Con(\mathrm{ZFC})$.  Its exact degree is $107286763576$ and its
circuit has $2204$ binary operations.  Its native base is $2^{683195}$
and its source-margin constant is
\[
 \boxed{\beta_S^{\mathrm{nz}}=2^{2743514}.}
\]
The four native ports together with this margin constant occupy
$685209$ raw bytes, compared with $1066835$ for the comparison
ten-integer instance.
\end{corollary}

This corollary uses the same imported K215, ordinary compiler/Pell,
and Sun theorems as the main construction.  The new work is the
fully specified composition and its coefficient and degree audit.

\subsection{The positive parent is compatible with the new tuple}

The literal parent polynomial is unchanged, including its old positive
transport coordinate $z_{\rm old}$.  To verify compatibility, use the
sixteen positive variables of the new retained59 system and write its
transport coordinate as $z_{\rm new}$.  Define
\begin{align*}
 q&=\beta J+1,& C&=q-F-\alpha,\\
 X&=wq,&Y&=sq^3,\\
 E&=XY,& k&=\eta+\zeta,\\
 c&=kY+\eta,& a&=E+Y,\\
 H&=4a+3,&\Delta&=a^2+H,\\
 D&=X+ac+\gamma H.
\end{align*}
The packed rank is
\[
 R=\bigl((q-1)(q-F)+\alpha\bigr)(q^2-1)
                  +(M_C+qM_F)J.
\]
Also put $V=of-c$, $I=ic^2$, $A_0=\Delta(f^2-1)$ and
$L_0=EkY$.  Denote the retained residuals by
\begin{align*}
 u_1&=\tau^2-L_0(L_0+k)-1,&
 u_2&=D^2-\Delta c^2-1,\\
 u_3&=A_0(V^2-y^2)+y^2-1,&
 u_4&=k-hE-R-1,\\
 u_5&=(K_0+w)C-F-z_{\rm new}(q-1),&
 u_6&=I^2-A_0,\\
 u_7&=jc-R-V.
\end{align*}
These are the seven equations already bound to the complete new tuple.

Extend a positive retained solution by
\begin{align*}
 z_{\rm old}&=z_{\rm new}+1,& N&=q^3,&L&=q^2-1,\\
 P&=E(c-\eta)+k,&C_2&=c^2,&F_2&=f^2,\\
 B_{\rm aux}&=I^2,&g&=y^2-V^2,
\end{align*}
and by the displayed expressions for the other parent auxiliaries.
The source coordinate $y$ is unchanged.  These parent coordinates
are distinct from the later ten-integer coordinates with the same
letter names.

\paragraph{The two sign obligations.}
All added auxiliaries are immediately positive except $V$ and $g$.
The inherited ordinary-kernel bootstrap gives
\[
 c>A\Delta^2,\qquad A=a+2,\qquad\Delta=A^2-1.
\]
This is equation~(3) of Section~4 in the byte-preserved source
\code{source_proofs/ORDINARY_BLANK_PROOF.md}.  It occurs before
auxiliary rank recovery: the first-norm index is at least $24$, the
main index is larger, and Pell growth gives the displayed bound.
In particular $c>2$ and $c^2>4\Delta$.  Here the ordinary strong
equation is
\[
 I^2=\Delta(f^2-1),\qquad I=ic^2.
\]
It therefore gives
\[
 f^2=1+\frac{i^2c^4}{\Delta}>4c^2.
\]
Consequently $V=of-c>c>1$, since $o$ is a positive integer.
The auxiliary equation now yields
\[
 I^2V^2-(I^2-1)y^2=1,\qquad
 g=y^2-V^2=\frac{V^2-1}{I^2-1}>0.
\]
These inequalities hold for every positive retained solution.

There is no identification of different normalized multipliers:
the parent uses the ordinary $i$.  In a canonical P67 witness,
$i_{\rm P67}=\psi_A(m)/c^2$, whereas the corresponding ordinary
parent coordinate is $i=\Delta i_{\rm P67}$.  The direct retained59
binding avoids conflating the two strong equations.

\paragraph{Exact residual identities.}
Under the displayed extension, all $22$ parent residuals vanish
identically except positions $8,12,13,14,15,21,22$.  In order, these
become
\[
 u_1,\quad u_2,\quad u_4-u_7,\quad u_7,\quad
 u_5,\quad u_6,\quad u_3+(V^2-y^2)u_6.
\]
The only packing simplification is the polynomial identity
\[
 (q-F)(N-q)-C(q^2-1)
 =\bigl((q-1)(q-F)+\alpha\bigr)(q^2-1)
 \qquad(N=q^3).
\]
The transport identity is
\[
 (K_0+w)C+q-F-z_{\rm old}(q-1)-1
 =(K_0+w)C-F-z_{\rm new}(q-1).
\]
Thus a positive retained solution extends to a positive literal
parent solution.

Conversely, a positive parent zero has $q>1$.  Its old transport
equation implies
\[
 (K_0+w)C=F+(z_{\rm old}-1)(q-1)>0,
\]
hence $C>0$ and $F<q$.  Restore the deleted packing register
\[
 G=q(q-F)-C=(q-1)C+q\alpha>0.
\]
The packing identity in the inherited quartic31 proof restores the
complete positive source.  Its ordinary kernel gives $X=wq=2^R$
and $R>3q$, so $w>q>F$.  Therefore
\[
 (z_{\rm old}-1)(q-1)=(K_0+w)C-F>0.
\]
It follows that $z_{\rm new}=z_{\rm old}-1$ is strictly positive.
The same identities recover the seven retained equations.  This
establishes the positive zero-set equivalence with all auxiliaries
and transport conventions accounted for.

The nonzero backend has already established the complete arithmetic
interface needed by this inherited kernel.  The positive parent
therefore has a zero precisely when there is a finite refutation.
Replacing its coordinates by $n_i+1$ gives the required source
quartic on $\N^{31}$.  Its value at zero remains positive because
the fifth residual is $1$ at the all-one parent assignment.

\subsection{Recomputing the source margin}

The new fixed port bit lengths are
\[
 (b_1,b_2,b_3,b_4)=(683195,688546,683195,683196).
\]
The symbolic coefficient recipes are unchanged.  For a weighted
coefficient written
$c_\alpha=\sum_e a_{\alpha,e}\prod_j A_j^{e_j}$, use the explicit
triangle bound
\[
 U_\alpha=\sum_e |a_{\alpha,e}|\,
                  2^{\sum_j e_jb_j}.
\]
Since $A_j<2^{b_j}$, $|c_\alpha|\le U_\alpha$.  All $402$ individual
bounds are recomputed from these recipes and the actual port
payloads.  Their largest certified exponent gives
\[
 |c_\alpha|+1<2^{2743487}.
\]
As before $|p_\alpha|\le|c_\alpha|$.  Because
$2\cdot24\cdot33^4=56924208<2^{27}$, the new even power
$\beta_S^{\rm nz}=2^{2743514}$ satisfies the full strict
source-margin hypothesis.  The new $\beta_S$ is materialized;
its raw length is $342940$ bytes.  Adding the four native payloads
gives $342269+342940=685209$ bytes, a $35.77\%$ reduction against
the five payloads of the comparison ten-integer construction.

The code weights, source interface, and outer relation-combining
formula now apply verbatim with the new values.  All required
properties have been re-established: a fixed quartic in $31$
natural variables, positive constant term, the exact factorial
coefficient convention, the $B_4$ code, and the strict even-power
margin.  No new existential parameter is introduced.

\subsection{Exact degree and reproducible composition}

The delivered arithmetic DAG is identical to the comparison DAG:
$485$ additions, $25$ subtractions and $1694$ multiplications,
with $10$ unknown leaves and $79$ fixed coefficient leaves.
The new fixed values cannot raise its degree bound.  To exclude
cancellation lowering it, the verifier computes the diagonal
leading coefficient with the actual new payloads:
\[
 \begin{array}{c|rrr}
 \text{prime}&101&1009&1000003\\ \hline
 \text{leading coefficient}&36&63&258712
 \end{array}
\]
Every entry is nonzero.  Thus
\[
 \deg F^{\rm nz}_{B_4}=113856T+27064=107286763576.
\]

The script \code{regenerate_binding.py} reads the sibling complete
backend and base parent, copies the four coherent payloads,
recomputes all coefficient bounds and $\beta_S$, and rewrites
all active binding metadata.  It is followed by
\code{b4_code.py}, \code{build_b4_sun10.py} and
\code{verify_b4_sun10.py}.  The independent review scripts
separately reconstruct the coefficients, check the byte binding,
and audit the formula and degree without importing these builders.
The source proofs and their Git blob identities are included.
The finite certificates supplement the universal transfer proof;
they do not replace the inherited arithmetic or logical theorems.


\section{Verification and proof status}
\label{sec:verification}
A zero-equivalence theorem is an assertion about all integer
assignments. Finite tests alone cannot establish it. The universal
new source argument and coefficient adapter are therefore proved
as ordinary mathematics, and the finite checks certify the
specific data and arithmetic claims used by those proofs.

For the main ten-integer equation, the executable verifier:
\begin{enumerate}
\item reconstructs the shifted sum of the 22 parent residual
squares and all 402 weighted coefficients exactly;
\item verifies the literal compressed ports, source hashes, and
all coefficient bounds;
\item checks all 52,360 four-multiset sums of the code;
\item expands the norm identities exactly in a formal polynomial
algebra, with 35 terms for three radicands and 495 terms for four;
\item compares each emitted circuit with the literal formulas
modulo three primes, including named intermediate quantities;
\item verifies the ten variable leaves, binary gate ledger,
fixed-leaf bindings, and an exact degree-attainment certificate.
\end{enumerate}
For the original-tuple comparison circuit, substituting a single indeterminate for
all unknowns gives a coefficient at degree $107286763576$ equal
to $16$ modulo 101. Since every gate also satisfies the advertised
degree upper bound, this nonzero residue proves exact degree.
The residues 599 modulo 1009 and 622356 modulo 1000003 provide
additional checks. One nonzero residue already suffices.
For the recommended composed instance the corresponding residues
are $36$, $63$, and $258712$. Its independent bound-and-binding
audit also reconstructs every weighted coefficient and checks
the new actual fixed payloads.

The native backend is checked independently without importing
the generator. The audit reconstructs all four integers and
checks 3,772 nonzero edge codes, 59,902 native-basis identities,
both rotation intersections, carry and padding bounds, and the
prime-power controller conditions on the actual new coefficients.
A separate mathematical review checks the order of the alignment
and occupancy arguments.

Clean rebuilding and byte comparison establish that the recorded
circuits are the generator outputs. They do not decide whether
the equation has a solution. The modular random assignments are
transcription diagnostics; the coefficient reconstruction, norm
identities, finite injectivity enumeration, and leading-coefficient
certificate have exact stated meanings.

\subsection{External and inherited dependencies}
The proof uses the pinned K215 and compiler equivalences, Sun's
2021 Lucas and binomial results, Sun's 2026 ten-integer bridge
and interpolation theorem, and elementary finite-field, Pell,
and unique-factorization facts. The 2026 paper remains a preprint
dependency. The new code and adapter have not been kernel-checked
end to end in Lean or Rocq.

The comparison ten-integer polynomial remains bound to the
original complete tuple. The recommended variant in
Section~\ref{composed:section} is fully rebound: all 402 weighted
coefficient recipes, their bounds, the actual source-base payload,
and the final degree certificate are regenerated from the new tuple.
Separate independent audits check both instances. The smaller tuple
is also concretely bound to P67 and the retained system.

\subsection{What the conclusions do and do not say}
The theorems are provable equivalences with consistency, not
proofs of consistency. The new degree reduction is established
relative to the explicitly pinned predecessor, and the
coefficient reduction is established relative to the complete
package-36 tuple. Neither is a claim that no smaller equation
exists. Classical Bose--Chowla sets and Sun's universal theorem
retain their attribution. The synthesis and explicit
instance improvements are the new contributions presented here;
a limited source search is not evidence of worldwide priority.

\section{Further research questions}
\label{sec:future}
\paragraph{Can the actual source support use a shorter additive code?}
A full $B_4$ code separates all 52,360 quartic monomials, while
the source has only 402 nonzero records. It is enough that each
supported exponent avoid every competing degree-at-most-four
monomial. An exact support-aware search could lower the maximum
position substantially. Checking collisions only between two
supported monomials would be insufficient: an unsupported
monomial from the fourth power could still contaminate the
coefficient being recovered.

\paragraph{Can span and circuit cost be optimized jointly?}
The former powers-of-five positions are unusually cheap to share.
The new irregular positions achieve a much shorter span but
increase the gate count. Codeword translation, assignment to
parent variables, and shared addition chains can be searched
together. A useful comparison must state both span and emitted
operation count.

\paragraph{Can the Jones nine-natural-variable route use additive coding?}
The related coefficient-extraction stage makes this plausible,
but each bound, mask, and reverse implication needs a separate
audit. The ten-integer proof does not automatically certify a
nine-variable circuit after a textual substitution. This is the
most direct route to combining the new degree improvement with
the collection's smallest natural-variable count.

\paragraph{Can other native backends benefit from the same adapter?}
Section~\ref{composed:section} completes the composition of the two
present improvements. The same architecture permits future coherent
compiler tuples, but each replacement must regenerate its source
coefficient bounds and fixed leaves. Automating that interface,
with proof certificates rather than only construction scripts,
would make subsequent backend comparisons much easier to audit.

\paragraph{Can the input quartic use fewer positive variables?}
The identity $D^2+c^2=(Ac)^2+1$ suggests removing a Pell
discriminant, but in the present system it merely moves the
auxiliary cost. Eliminating the index quotient runs into
divisibility and proof-order obstacles involving
$\gcd(E,H)$. In particular $X=2^R$ cannot be used before the
argument establishing the packed index has recovered it.
No reduced quartic count or eleven-positive-variable construction
is asserted here.

\paragraph{Which formal lemma should come first?}
The most reusable addition to the formal development is the
additive coefficient-extraction lemma parameterized by a finite
injectivity certificate, followed by the mask and source
interface. This separates the new proof from the existing
Lucas and machine endpoints. An exact circuit checker can
verify degree upper bounds and the nonzero modular
leading coefficient.

\paragraph{How close is the displayed code to a minimum?}
For an integer $B_4$ code with 32 positions in $[0,T]$,
counting gives $4T+1\ge\binom{35}{4}$, hence $T\ge13090$.
This lower bound is weak and does not show that 942302 is close
to optimal. Optimizing this finite code is a concrete
combinatorial task. Proving a globally simplest Diophantine
consistency equation is a separate problem.

\paragraph{Can occupancy information be compressed further?}
The new per-edge widths are minimal under the stated separate
injective nonzero-word interface. This leaves open joint
encodings, different tile presentations, alternative support
geometry, and larger changes to the computation-to-polynomial
compiler. A claimed gain must preserve both the local language
and the global alignment/halting theorem.

\section{Reproduction and incorporation}
The top-level \code{reproduce.py} runs the construction and
finite audits using Python's standard library. The report is
built with \code{latexmk -pdf report.tex}. The directory includes:
\begin{itemize}
\item the complete LaTeX article and its compiled PDF;
\item the composed ten-integer circuit, the original-tuple comparison
instance, and four alternative finalizers for that comparison;
\item the literal parent, coefficient recipes, and compressed
fixed integers;
\item the new native backend, its concrete P67/system bindings,
and independent coefficient audit;
\item finite verification receipts, independent mathematical
reviews, source pins, notices, and licenses.
\end{itemize}
The package is intended as a new sibling in the user's
\code{con-zfc-diophantine-forms} collection. Existing source
numbers and old tuples should retain their identities.
The supplied fixed code positions describe a new polynomial;
they are not an edit to the old package-31 polynomial's degree.

\appendix
\section{The literal positive parent system}
\label{app:parent}
Use the variables and affine abbreviations of
Section~\ref{sec:repo}, and write $R=jc-V$ as a computed
expression. The 22 equations below specify the parent of the
two ten-integer instances. Their residuals are left side minus right
side.
\begin{align*}
 L&=q^2-1,&N&=q(L+1),\\
 X&=wq,&a&=E+sN,\\
 E&=X(a-E),&c&=k(a-E)+\eta,\\
 P&=E(c-\eta)+k,&\tau^2&=P(P-k)+1,\\
 D&=X+ac+\gamma H,&\Delta&=a^2+H,\\
 C_2&=c^2,&D^2&=\Delta C_2+1,\\
 k&=R+1+hE.
\end{align*}
The next two transport equations are
\begin{align*}
 R&=(q-F)(N-q)-CL+(M_C+qM_F)J,\\
 (K_0+w)C+(q-F)-z(q-1)&=1.
\end{align*}
Finally,
\begin{align*}
 F_2&=f^2,&V&=of-c,\\
 g&=y^2-V^2,&I&=iC_2,\\
 B_{\rm aux}&=I^2,&B_{\rm aux}&=\Delta(F_2-1),\\
 y^2-B_{\rm aux}g&=1.
\end{align*}
The transport coordinate here is the original unshifted one.
The later 60-operation quadratic system uses a shifted transport
coordinate; its residual list must not be silently substituted
when reconstructing this package's coefficient records.
The JSON source is authoritative and is checked against the
displayed equations.

\begin{thebibliography}{99}
\bibitem{Repo}
V.~Reshetnikov, \emph{ProveIt}, revision
\code{22e75dd6b87b1b199bb08c778d7cf075ae5aae08},
23 October 2026. Collection:
\url{https://github.com/VladimirReshetnikov/ProveIt/tree/22e75dd6b87b1b199bb08c778d7cf075ae5aae08/SetTheory/Cardinals/docs/reports/hilbert-tenth-problem/con-zfc-diophantine-forms}.
The source provenance ledger identifies individual paths and blobs.

\bibitem{Package31}
ProveIt package \code{31-sun10}, \emph{Smaller explicit Diophantine
forms for Con(ZFC)}, at revision \cite{Repo}, particularly
\code{signed10/SIGNED10_PROOF.md},
\code{coefficients/COEFFICIENT_PROOF.md}, and
\code{signed10/data/parent_quartic31.json}.
Its ten-integer finalizer improvements are inherited here.

\bibitem{Package36}
ProveIt package \code{36-small-tuple},
\emph{Smaller Diophantine presentations of ZFC consistency},
at revision \cite{Repo}. The native compiler, actual-rank
padding method, and 59-operation system are the baseline
for the separate coefficient improvement.

\bibitem{Sun2021}
Z.-W.~Sun, \emph{Further results on Hilbert's Tenth Problem},
Science China Mathematics \textbf{64} (2021), 281--306.
\url{https://arxiv.org/abs/1704.03504v7}.
The source coding and Lucas/binomial results are used with
explicit hypotheses.

\bibitem{Sun2026}
Z.-W.~Sun, \emph{Ten unknowns for Hilbert's tenth problem over the
integers}, arXiv:2609.01594v2, 2 September 2026.
\url{https://arxiv.org/abs/2609.01594v2}.
Proposition 2.1, Theorems 3.1 and 4.1, Propositions 6.1--6.2,
and Section 7. This is a preprint dependency.

\bibitem{BoseChowla}
R.~C.~Bose and S.~Chowla, \emph{Theorems in the additive theory of
numbers}, Commentarii Mathematici Helvetici \textbf{37}
(1962/63), 141--147.
\url{https://doi.org/10.1007/BF02566968}.
Earlier author report:
\url{https://www.cs.umd.edu/~gasarch/COURSES/858/S13/BoseChowla.pdf}.

\bibitem{Bayer2025}
J.~Bayer, M.~David, M.~Ha{\ss}ler, Y.~Matiyasevich, and
D.~Schleicher, \emph{Diophantine equations over $\Z$: universal
bounds and parallel formalization}, arXiv:2506.20909v3,
22 October 2026; first version 2025.
\url{https://arxiv.org/abs/2506.20909v3}.
Equations (1.6)--(1.7) and Definition 2.1 give an independent
exposition of the positional source architecture.
\end{thebibliography}
\end{document}

